\documentclass[12pt,a4paper]{article}

% ── Packages ─────────────────────────────────────────────────────────────────
\usepackage[T1]{fontenc}
\usepackage[utf8]{inputenc}
\usepackage{lmodern}
\usepackage[margin=2.5cm]{geometry}
\usepackage{amsmath,amssymb}
\usepackage{graphicx}
\usepackage{caption}
\usepackage{subcaption}
\usepackage{booktabs}
\usepackage{hyperref}
\usepackage{cite}
\usepackage{setspace}
\usepackage{microtype}
\usepackage{xcolor}
\usepackage{parskip}

\hypersetup{
  colorlinks = true,
  linkcolor  = blue!60!black,
  citecolor  = green!50!black,
  urlcolor   = blue!60!black,
}

% ── Title block ───────────────────────────────────────────────────────────────
\title{%
  \textbf{Influence of Hydrogen Enrichment of Methane on
  Laminar Flame Speed and NO$_x$ Emissions}\\[0.8em]
  \large Computational Methods in Combustion
}
\author{Jędrzej Koc\\[0.3em]
}
\date{\today}

% =============================================================================
\begin{document}
\maketitle
\thispagestyle{empty}

% ─────────────────────────────────────────────────────────────────────────────


\newpage
\tableofcontents
\newpage

% =============================================================================
\section{Introduction}
\label{sec:intro}

The global energy sector is under increasing pressure to reduce greenhouse gas
emissions and transition toward sustainable energy carriers. Natural gas, composed
predominantly of methane, currently plays a central role in power generation and
industrial heating owing to its relatively low carbon intensity compared with coal
or oil. However, even complete combustion of methane releases CO$_2$, and the
transition to a truly low-carbon energy system requires either carbon capture or a
shift toward carbon-free fuels.

Hydrogen is widely regarded as a key element of this transition. It releases no
CO$_2$ upon combustion, has a very high energy content per unit mass, and can be
produced from renewable electricity via electrolysis (so-called green hydrogen).
A pragmatic near-term pathway to decarbonisation is the blending of H$_2$ into
the existing natural gas grid---a strategy already tested at the distribution
scale in several countries. Such blending enables incremental emission reductions
while the infrastructure, safety standards, and economic frameworks for pure
hydrogen supply chains mature.

From a combustion-engineering perspective, however, hydrogen addition to methane
significantly alters flame behaviour. The high diffusivity and low ignition energy
of H$_2$ can improve flame stability, but also raise safety concerns related to
flashback and detonation risk. Changes in laminar flame speed, radical chemistry,
and adiabatic flame temperature all directly affect burner design, emission
formation, and operability limits. A quantitative, mechanistically grounded
understanding of these effects is therefore essential for engineers tasked with
adapting existing combustion equipment to hydrogen-enriched fuels.

This report addresses this need through a computational parametric study. The
laminar flame speed $S_u$, maximum flame temperature $T_\mathrm{max}$, NO
emissions, radical peak concentrations, and spatial flame structure are evaluated
as functions of hydrogen volume fraction in the CH$_4$/H$_2$ fuel blend.

% =============================================================================
\section{Literature Review}
\label{sec:lit}

The combustion characteristics of hydrogen-enriched natural gas (HENG) have been
the subject of extensive research over the past two decades. The following
paragraphs summarise the principal findings relevant to the present study.

\textbf{Laminar flame speed.} It is well established that the laminar burning
velocity of CH$_4$/H$_2$/air mixtures increases monotonically with hydrogen
fraction. Huang~et~al.~\cite{Huang2006} and Halter~et~al.~\cite{Halter2005}
measured this increase experimentally and attributed it to both the significantly
higher intrinsic burning velocity of pure hydrogen ($\approx 200\,\mathrm{cm/s}$
at stoichiometry) and to enhanced chain-branching kinetics that H$_2$ introduces
into the methane oxidation pathway. Di~Sarli and Di~Benedetto~\cite{DiSarli2007}
confirmed through numerical simulation that even moderate H$_2$ additions (below
30\,\%) produce a disproportionate increase in $S_u$ relative to what a linear
blending rule would predict.

\textbf{Radical pool and chemical kinetic enhancement.} The kinetic
acceleration of HENG flames is closely linked to the enrichment of the radical
pool, particularly OH, H, and O. Hydrogen addition provides additional H atoms
via thermal decomposition, which then participate in the key chain-branching
reaction H + O$_2$ $\rightarrow$ OH + O. This positive feedback loop greatly
amplifies the overall reactivity of the mixture beyond the contribution of H$_2$
as a pure energy source. Studies using sensitivity analysis
\cite{Schefer2002,Coppens2007} consistently identify the H + O$_2$ reaction as
the most sensitive step, and confirm that radical-pool enhancement is the dominant
kinetic mechanism at low-to-moderate H$_2$ blending ratios.

\textbf{Adiabatic flame temperature and NOx.} Theoretical considerations and
measurements agree that H$_2$ addition raises the adiabatic flame temperature at
stoichiometric conditions, primarily because the lower carbon-to-hydrogen ratio
reduces heat capacity while the combustion enthalpy per unit volume remains high.
This temperature increase has a strong positive effect on thermal NO formation via
the extended Zeldovich mechanism \cite{Zeldovich1947}, and numerous studies
\cite{Guo2010,Zhang2007} report elevated NO$_x$ emissions in HENG flames
proportional to the temperature increase. The consensus is that the NOx penalty
is a significant engineering challenge that must be addressed through staged
combustion, lean premixing, or exhaust gas recirculation.

\textbf{Flame structure.} Spatially resolved measurements and simulations show
that hydrogen enrichment decreases the flame thickness---defined as the width of
the preheat zone---due to the higher thermal diffusivity of H$_2$ and the more
compact reaction zone. This results in a narrower, more intense heat-release
profile and influences flame--turbulence interaction in practical burners.

% =============================================================================
\section{Model Description}
\label{sec:model}

\subsection{Software and Kinetic Mechanism}

All simulations were performed using Cantera~3.2 \cite{Cantera}, an open-source
object-oriented software package for chemical kinetics and reacting-flow
calculations. The one-dimensional freely propagating premixed flame model
(\texttt{ct.FreeFlame}) was employed, which solves the steady-state conservation
equations for mass, species, and energy on an adaptive grid using the Newton
method with automatic Jacobian evaluation.

The GRI~Mech~3.0 mechanism \cite{Smith1999} (\texttt{gri30.yaml}) was used for
all simulations. It contains 53 species and 325 elementary reactions, covering
methane oxidation, C$_2$ chemistry, and the full NO$_x$ formation sub-mechanism
including the thermal (Zeldovich), prompt (Fenimore), and N$_2$O pathways. This
makes it the standard choice for studying NOx formation in methane-based premixed
flames.

\subsection{Flame Configuration and Boundary Conditions}

The computational setup is summarised in Table~\ref{tab:conditions}.

\begin{table}[htbp]
  \centering
  \caption{Simulation parameters.}
  \label{tab:conditions}
  \begin{tabular}{lll}
    \toprule
    Parameter & Symbol & Value \\
    \midrule
    Unburned gas temperature & $T_u$ & 300\,K \\
    Pressure                 & $P$   & 1\,atm (101\,325\,Pa) \\
    Equivalence ratio        & $\phi$ & 1.0 \\
    H$_2$ fraction in fuel   & $x_{\mathrm{H_2}}$ & 0, 10, 20, 30, 40, 50\,vol.\% \\
    CH$_4$ fraction in fuel  & $x_{\mathrm{CH_4}}$ & $1 - x_{\mathrm{H_2}}$ \\
    Oxidiser                 & ---   & Dry air (21\,\% O$_2$, 79\,\% N$_2$) \\
    \bottomrule
  \end{tabular}
\end{table}

\subsection{Mixture Composition}

For each hydrogen fraction $x_{\mathrm{H_2}}$ (mole fraction of H$_2$ in the
fuel blend), the stoichiometric oxygen requirement per mole of fuel is:

\begin{equation}
  n_{\mathrm{O_2,stoich}} = 2\,x_{\mathrm{CH_4}} + 0.5\,x_{\mathrm{H_2}}
  \label{eq:O2_stoich}
\end{equation}

where $x_{\mathrm{CH_4}} = 1 - x_{\mathrm{H_2}}$. At equivalence ratio $\phi$,
the actual oxygen supplied is $n_{\mathrm{O_2}} = n_{\mathrm{O_2,stoich}}/\phi$,
and the nitrogen is set by the air composition:
$n_{\mathrm{N_2}} = n_{\mathrm{O_2}} \times (79/21)$.

\subsection{Grid and Solver Settings}

The initial domain width was set to 3\,cm. Automatic grid refinement was enabled
with the following criteria: ratio~3, slope~0.06, curve~0.12, and a prune
threshold of 0.01. These settings ensure that steep gradients in temperature and
species profiles are adequately resolved while avoiding unnecessary grid points in
smooth regions. The solver was allowed to iterate through three stages of
increasing refinement until all criteria were satisfied (\texttt{auto=True}).

% =============================================================================
\section{Results and Discussion}
\label{sec:results}

\subsection{Laminar Flame Speed}

\begin{figure}[htbp]
  \centering
  \includegraphics[width=0.75\textwidth]{flame_speed.png}
  \caption{Laminar flame speed $S_u$ as a function of hydrogen fraction in the
           CH$_4$/H$_2$ fuel blend at $\phi = 1.0$, $T_u = 300\,\mathrm{K}$,
           $P = 1\,\mathrm{atm}$.}
  \label{fig:flame_speed}
\end{figure}

Figure~\ref{fig:flame_speed} shows a clear, super-linear increase in the laminar
flame speed with hydrogen enrichment. The pure methane flame ($x_{\mathrm{H_2}}
= 0$) attains a speed of approximately 37--38\,cm/s, consistent with the
well-established experimental reference value for CH$_4$/air at stoichiometry.
As H$_2$ content rises to 50\,vol.\%, the flame speed increases by more than
50\,\%, reaching values above 55\,cm/s.

Three physical mechanisms contribute to this acceleration. First, H$_2$ itself
has an intrinsically high burning velocity ($\approx\!200\,\mathrm{cm/s}$), and
even small additions shift the mixture's effective reactivity. Second, hydrogen's
high thermal diffusivity ($\alpha_{\mathrm{H_2}} \approx 5 \times
\alpha_{\mathrm{CH_4}}$) enhances preheat-zone conduction, allowing the unburned
gas to reach ignition temperature more rapidly. Third, and most importantly at
moderate blending ratios, H$_2$ dramatically amplifies the radical pool
(discussed in Section~\ref{sec:radicals}), accelerating the rate-limiting
chain-branching steps of hydrocarbon oxidation.

The super-linear character of the increase---visible in Fig.~\ref{fig:flame_speed}
as an upward curvature---is a well-known feature of HENG flames and is explained
by the synergistic interaction between H$_2$ kinetics and the methane oxidation
network rather than a simple mixing effect.

\subsection{Maximum Flame Temperature}

\begin{figure}[htbp]
  \centering
  \includegraphics[width=0.75\textwidth]{flame_temperature.png}
  \caption{Maximum adiabatic flame temperature as a function of hydrogen fraction.}
  \label{fig:T_max}
\end{figure}

The maximum flame temperature shown in Figure~\ref{fig:T_max} increases
approximately linearly with hydrogen content, rising by roughly 80\,K over the
studied range. This increase arises from two thermodynamic effects. First,
replacing part of the CH$_4$ (which produces CO$_2$ and H$_2$O upon combustion)
with H$_2$ (which produces only H$_2$O) reduces the specific heat of the product
mixture, since CO$_2$ has a higher molar heat capacity than H$_2$O. Second, the
higher hydrogen-to-carbon ratio means proportionally more water vapour in the
products, and the combustion enthalpy released per unit mole of oxygen consumed
is slightly higher for H$_2$ than for CH$_4$. Both effects result in a higher
adiabatic flame temperature, which has direct consequences for thermal NO
formation (see Section~\ref{sec:NOx}).

\subsection{NO Emissions}
\label{sec:NOx}

\begin{figure}[htbp]
  \centering
  \includegraphics[width=0.75\textwidth]{NO_emission.png}
  \caption{NO mole fraction (as ppm) in the burned gas as a function of hydrogen
           fraction.}
  \label{fig:NO}
\end{figure}

Figure~\ref{fig:NO} demonstrates a significant increase in NO emissions with
hydrogen enrichment. The dominant formation pathway in premixed methane/air flames
near stoichiometry is the thermal (Zeldovich) mechanism:

\begin{align}
  \mathrm{N_2} + \mathrm{O}  &\;\rightleftharpoons\; \mathrm{NO} + \mathrm{N} \label{eq:Z1}\\
  \mathrm{N}   + \mathrm{O_2} &\;\rightleftharpoons\; \mathrm{NO} + \mathrm{O} \label{eq:Z2}\\
  \mathrm{N}   + \mathrm{OH}  &\;\rightleftharpoons\; \mathrm{NO} + \mathrm{H} \label{eq:Z3}
\end{align}

The rate of Reaction~\eqref{eq:Z1} has an extremely strong temperature dependence
(activation energy $\approx\!314\,\mathrm{kJ/mol}$), expressed by:

\begin{equation}
  k_1 = A\,T^n \exp\!\left(-\frac{E_a}{RT}\right)
  \label{eq:Arrhenius}
\end{equation}

The 80\,K increase in $T_\mathrm{max}$ shown in Fig.~\ref{fig:T_max} therefore
translates into a disproportionately large increase in the Zeldovich NO formation
rate, explaining the steep rise in NO concentration seen in Fig.~\ref{fig:NO}.
The elevated OH radical concentration (Fig.~\ref{fig:radicals}) further promotes
Reaction~\eqref{eq:Z3}, creating an additional feedback loop between radical
enrichment and NOx formation.

This result underscores a fundamental trade-off in HENG combustion: hydrogen
addition improves burning velocity and combustion stability, but simultaneously
increases the thermal NOx penalty. Mitigation strategies---such as lean premixing,
exhaust gas recirculation, or selective catalytic reduction (SCR)---are therefore
critical for practical HENG burner applications.

\subsection{Radical Pool Enhancement}
\label{sec:radicals}

\begin{figure}[htbp]
  \centering
  \includegraphics[width=0.75\textwidth]{radicals.png}
  \caption{Peak mole fractions of OH, H, and O radicals as a function of hydrogen
           fraction.}
  \label{fig:radicals}
\end{figure}

Figure~\ref{fig:radicals} reveals that the peak concentrations of all three key
flame radicals---OH, H, and O---increase monotonically with hydrogen enrichment.
The H radical, in particular, shows a pronounced increase, consistent with H$_2$
decomposition acting as an additional source of H atoms in the preheat zone via:

\begin{equation}
  \mathrm{H_2} + \mathrm{M} \;\rightleftharpoons\; \mathrm{H} + \mathrm{H} + \mathrm{M}
\end{equation}

and via abstraction by O:
\begin{equation}
  \mathrm{H_2} + \mathrm{O} \;\rightleftharpoons\; \mathrm{OH} + \mathrm{H}
\end{equation}

The resulting surplus of H atoms accelerates the critical chain-branching step:
\begin{equation}
  \mathrm{H} + \mathrm{O_2} \;\rightarrow\; \mathrm{OH} + \mathrm{O}
  \label{eq:chainbranch}
\end{equation}

which simultaneously increases OH and O concentrations. The expanded radical pool
accelerates the overall oxidation rate of both CH$_4$ and CO, leading to the
higher laminar flame speeds observed in Fig.~\ref{fig:flame_speed}. This
mechanism also explains the super-linear dependence of $S_u$ on H$_2$ content:
the radical chain-branching creates an autocatalytic amplification effect that
grows faster than the linear increase in H$_2$ concentration alone would suggest.

\subsection{Spatial Flame Structure}

\begin{figure}[htbp]
  \centering
  \includegraphics[width=0.75\textwidth]{temperature_profile.png}
  \caption{Spatial temperature profiles for the 0\,\% and 50\,\% H$_2$ cases
           (grids shifted so the point of maximum temperature gradient is at
           $x = 0$).}
  \label{fig:T_profile}
\end{figure}

\begin{figure}[htbp]
  \centering
  \includegraphics[width=0.75\textwidth]{HRR_profile.png}
  \caption{Spatial heat release rate (HRR) profiles for the 0\,\% and 50\,\%
           H$_2$ cases.}
  \label{fig:HRR}
\end{figure}

Figures~\ref{fig:T_profile} and~\ref{fig:HRR} compare the spatial flame
structure for the two extreme cases: pure methane and 50\,\% hydrogen enrichment.
Both grids have been shifted so that the point of maximum temperature gradient
coincides with $x = 0$, enabling a direct comparison of flame structure rather
than absolute flame position.

\textbf{Temperature profile (Fig.~\ref{fig:T_profile}).} The 50\,\% H$_2$ flame
exhibits a steeper temperature rise over a shorter spatial distance. This narrowing
of the preheat zone is a direct consequence of the higher thermal diffusivity of
H$_2$ relative to CH$_4$: a larger fraction of the heat released in the reaction
zone is conducted upstream into the unburned mixture, raising the local temperature
gradient and, by Le Chatelier's principle applied to the energy equation,
compressing the preheat zone. The final burned-gas temperature is also higher, as
expected from Fig.~\ref{fig:T_max}.

\textbf{Heat release rate profile (Fig.~\ref{fig:HRR}).} The peak HRR is
significantly higher for the 50\,\% H$_2$ case, and the reaction zone (defined
by the region where HRR $> 10$\,\% of peak) is markedly narrower. This behaviour
is consistent with the higher laminar flame speed: since the flame consumes the
reactant mixture faster per unit area, the volumetric heat release rate must
increase to supply the required power output. The narrowing of the HRR profile
has important implications for turbulent combustion: a thinner reaction zone is
more susceptible to turbulent stretching and local quenching, potentially
affecting the lean blowout limit at elevated hydrogen fractions.

% =============================================================================
\section{Conclusions}
\label{sec:conclusions}

The following key engineering conclusions are drawn from this computational study
of hydrogen-enriched methane combustion:

\begin{itemize}
  \item \textbf{Laminar flame speed increases substantially with H$_2$ addition.}
        From $\approx\!37\,\mathrm{cm/s}$ for pure CH$_4$, the speed rises to
        $>\!55\,\mathrm{cm/s}$ at 50\,\% H$_2$, a gain of over 50\,\% driven by
        radical-pool kinetics and enhanced thermal diffusivity.

  \item \textbf{Maximum flame temperature rises by $\approx\!80$\,K.}
        The thermodynamic shift caused by replacing CO$_2$-producing methane
        combustion with H$_2$O-producing hydrogen combustion increases the
        adiabatic flame temperature monotonically.

  \item \textbf{NO emissions increase significantly.}
        The Zeldovich thermal mechanism, highly sensitive to temperature, produces
        substantially more NO as the flame temperature rises, presenting a critical
        challenge for emission compliance.

  \item \textbf{Radical pool (OH, H, O) is enriched by H$_2$.}
        The chain-branching cycle H + O$_2$ $\rightarrow$ OH + O is strongly
        amplified by the additional H atoms provided by H$_2$ decomposition,
        explaining the super-linear increase in flame speed.

  \item \textbf{The reaction zone narrows and intensifies with H$_2$ addition.}
        Spatial profiles show a thinner preheat zone and a higher, narrower HRR
        peak, implying increased susceptibility to turbulent disruption and
        flashback at high H$_2$ fractions.

  \item \textbf{Hydrogen blending is a viable but constrained decarbonisation
        strategy.}
        While HENG reduces CO$_2$ emissions proportionally to the hydrogen
        fraction, the NOx and flashback penalties require active mitigation
        through burner redesign, lean premixing, or post-combustion treatment.
\end{itemize}

% =============================================================================

\begin{thebibliography}{99}

\bibitem{Huang2006}
  Z.~Huang, Y.~Zhang, K.~Zeng, B.~Liu, Q.~Wang, D.~Jiang,
  \textit{Measurements of laminar burning velocities for natural gas-hydrogen-air
  mixtures},
  Combustion and Flame 146 (2006) 302--311.

\bibitem{Halter2005}
  F.~Halter, C.~Chauveau, N.~Djebaili-Chaumeix, I.~Gökalp,
  \textit{Characterization of the effects of pressure and hydrogen concentration
  on laminar burning velocities of methane--hydrogen--air mixtures},
  Proceedings of the Combustion Institute 30 (2005) 201--208.

\bibitem{DiSarli2007}
  V.~Di~Sarli, A.~Di~Benedetto,
  \textit{Laminar burning velocity of hydrogen--methane/air premixed flames},
  International Journal of Hydrogen Energy 32 (2007) 637--646.

\bibitem{Schefer2002}
  R.~W.~Schefer,
  \textit{Hydrogen enrichment for improved lean flame stability},
  International Journal of Hydrogen Energy 28 (2003) 1131--1141.

\bibitem{Coppens2007}
  F.~H.~V.~Coppens, J.~De~Ruyck, A.~A.~Konnov,
  \textit{The effects of composition on burning velocity and nitric oxide
  formation in laminar premixed flames of CH$_4$ + H$_2$ + O$_2$ + N$_2$},
  Combustion and Flame 149 (2007) 409--417.

\bibitem{Zeldovich1947}
  Ya.~B.~Zel'dovich,
  \textit{The oxidation of nitrogen in combustion and explosions},
  Acta Physicochimica URSS 21 (1946) 577--628.

\bibitem{Guo2010}
  H.~Guo, B.~Smallwood, F.~Liu, Y.~Ju, Ö.~L.~Gülder,
  \textit{The effect of hydrogen addition on flammability limit and NOx emission
  in ultra-lean counterflow CH$_4$/air premixed flames},
  Proceedings of the Combustion Institute 30 (2005) 303--311.

\bibitem{Zhang2007}
  Y.~Zhang, Z.~Huang, L.~Wei, J.~Zhang, C.~K.~Law,
  \textit{Experimental and modeling study on ignition delays of lean mixtures of
  methane, hydrogen, oxygen, and argon at elevated pressures},
  Combustion and Flame 150 (2007) 1--13.

\bibitem{Cantera}
  D.~G.~Goodwin, R.~L.~Speth, H.~K.~Moffat, B.~W.~Weber,
  \textit{Cantera: An object-oriented software toolkit for chemical kinetics,
  thermodynamics, and transport processes},
  Version 3.2, 2024, \url{https://www.cantera.org}.

\bibitem{Smith1999}
  G.~P.~Smith, D.~M.~Golden, M.~Frenklach, N.~W.~Moriarty, B.~Eiteneer,
  M.~Goldenberg, C.~T.~Bowman, R.~K.~Hanson, S.~Song, W.~C.~Gardiner,
  V.~V.~Lissianski, Z.~Qin,
  GRI-Mech 3.0, 1999, \url{http://combustion.berkeley.edu/gri-mech/}.

\end{thebibliography}

\end{document}
